Ce thème suit une seule question : \emph{de quoi parle-t-on quand on dit « nombre » ?} La
réponse change quatre fois au cours de la scolarité, et chaque changement est motivé par un
calcul qu'on ne savait pas faire.

On part des entiers et de la division euclidienne, dont tout le collège découle : parité,
critères de divisibilité, pgcd, nombres premiers, décomposition. Vient le besoin de partager,
donc les fractions, puis celui de mesurer une diagonale, donc les irrationnels — et
$\sqrt{2}$, dont l'irrationalité est le premier théorème du programme qui affirme une
\emph{impossibilité}. Le lycée reprend ces objets avec un outillage neuf : valeur absolue,
puissances, sommes, récurrence. Enfin, l'impossibilité de résoudre $x^2 = -1$ ouvre les
complexes, où l'équation du second degré a toujours deux racines.

Le thème se clôt sur l'arithmétique de la spécialité — congruences, Bézout, Gauss, petit
théorème de Fermat — qui est exactement la divisibilité de sixième, reprise avec les moyens
de la terminale. Les deux chapitres se lisent l'un après l'autre sans rupture, et c'est
précisément la raison de les avoir mis dans le même thème.
